%% =====================================================================
%% FeatureLiftBench -- full paper draft
%% ACM acmart / anonymous review version
%%
%% IMPORTANT:
%%   1) Keep anonymous mode for double-anonymous review.
%%   2) Replace every \draftnote{...} before submission.
%%   3) Headline Python-150 results are populated from the frozen final analysis.
%% =====================================================================

\documentclass[acmsmall,screen,review,anonymous]{acmart}

%TC:macro \cite [option:text,text]
%TC:macro \citep [option:text,text]
%TC:macro \citet [option:text,text]
%TC:envir table 0 1
%TC:envir table* 0 1
%TC:envir tabular [ignore] word
%TC:envir displaymath 0 word
%TC:envir math 0 word
%TC:envir comment 0 0

\usepackage{booktabs}
\usepackage{tabularx}
\usepackage{listings}
\usepackage{xspace}

\newcommand{\flb}{\textsc{FeatureLiftBench}\xspace}
\newcommand{\rres}{\textsc{RRES}\xspace}
\newcommand{\mainprotocol}{\textsc{Main}\xspace}
\newcommand{\validatoragent}{\textsc{Validator-Agent}\xspace}
\newcommand{\draftnote}[1]{\textcolor{red}{[TODO: #1]}}


%% =====================================================================
%% FROZEN EXPERIMENTAL RECORD -- Python-150 freeze v2 (2026-09-04/05)
%% =====================================================================
%% SOURCE OF TRUTH FOR ALL NUMBERS BELOW:
%%   reports/paper_analysis/python150_prime_v2_analysis_20260905/task_results.csv
%% REPRODUCTION:
%%   .venv/bin/python reports/paper_analysis/python150_paper_analysis_final/build.py
%% ACTIVE FREEZE ID:
%%   6c20ff0307762503a73cbb9ff32e9992c6446e4b17483a68373027be58cbf419
%% PROTOCOL:
%%   OpenHands Official Main; one attempt/task; 120 steps; 128k context;
%%   Full-Repository / No-Hint; timeout 3600s; empty submission = failure;
%%   Functional = Build AND Public AND Hidden AND Isolation; run.status is NOT score.
%% RELEASE / HEADLINE SPLITS:
%%   Released Python suite: 200 tasks / 176 repos.
%%   Headline frozen Python-150: 150 tasks / 127 repos.
%%   Core-100 + hard3-50 are the in-suite construction split.
%%   Official additional 50-task release split is appendix-only and is NOT the
%%   paper's main difficulty claim.
%% MODELS (headline Python-150):
%%   Pro   = DeepSeek V4 Pro
%%   Flash = DeepSeek V4 Flash
%%   Luna  = GPT-5.6 Luna (OpenLux)
%%   GLM   = GLM-5.3-Flash
%%   Qwen  = Qwen3.6-35B
%%   OSS   = GPT-OSS 120B
%%
%% MAIN FUNCTIONAL PASS TABLE:
%%   Pro   115/150 = 76.7%, Wilson 69.3--82.7%; Core 94/100; hard3 21/50; empty 0
%%   Flash 108/150 = 72.0%, Wilson 64.3--78.6%; Core 90/100; hard3 18/50; empty 0
%%   Luna  102/150 = 68.0%, Wilson 60.2--74.9%; Core 82/100; hard3 20/50; empty 6
%%   GLM    68/150 = 45.3%, Wilson 37.6--53.3%; Core 62/100; hard3  6/50; empty 38
%%   Qwen   63/150 = 42.0%, Wilson 34.4--50.0%; Core 55/100; hard3  8/50; empty 25
%%   OSS    36/150 = 24.0%, Wilson 17.9--31.4%; Core 25/100; hard3 11/50; empty 2
%%   Unsolved by all six = 28/150; of these, hard3 = 22.
%%
%% EXACT McNEMAR (two-sided exact binomial):
%%   Pro vs Flash: Pro-only 10, Flash-only 3, p=0.09229  [DO NOT claim significant]
%%   Pro vs Luna : Pro-only 18, Luna-only  5, p=0.01062
%%   GLM vs Qwen : GLM-only 27, Qwen-only 22, p=0.5682   [DO NOT claim GLM > Qwen]
%%
%% MUTUALLY EXCLUSIVE FIRST OUTCOME [Pass, Missing, Build, Public, Hidden, Isolation]:
%%   Pro   [115, 0,  0, 25, 10, 0]
%%   Flash [108, 0,  0, 26, 15, 1]
%%   Luna  [102, 6,  3, 27, 12, 0]
%%   GLM   [ 68,38,  3, 31,  8, 2]
%%   Qwen  [ 63,25,  6, 35, 21, 0]
%%   OSS   [ 36, 2, 18, 70, 23, 1]
%%
%% NON-EXCLUSIVE GATES ON DELIVERED PACKAGES
%% [has_package, build_fail, public_fail, hidden_fail, isolation_fail,
%%  isolation_given_build, isolation_residual]:
%%   Pro   [150, 0,25, 29, 0,0,0]
%%   Flash [150, 0,26, 36, 1,1,1]
%%   Luna  [144, 3,30, 35, 3,0,0]
%%   GLM   [112, 3,34, 38, 5,2,2]
%%   Qwen  [125, 6,41, 58, 8,2,0]
%%   OSS   [148,18,88,100,22,4,1]
%%   Raw Isolation failures total = 39; binding Isolation residual total = 4.
%%   Safe claim: Independence is rarely the binding constraint; preserving
%%   complete behavior is the dominant residual challenge.
%%
%% EMPIRICAL SOLVE-FREQUENCY SPECTRUM [Core, hard3, total]:
%%   0/6 [ 6,22,28]
%%   1/6 [ 0, 7, 7]
%%   2/6 [ 4, 5, 9]
%%   3/6 [17, 5,22]
%%   4/6 [31, 5,36]
%%   5/6 [27, 4,31]
%%   6/6 [15, 2,17]
%%
%% LIFT TYPE x SPLIT COUNTS:
%%   Core-100: Direct 53, Adapted 44, Composite 3
%%   hard3-50: Direct 3, Adapted 32, Composite 15
%%   Therefore raw Direct>Adapted>Composite is confounded.
%% LOGISTIC: passed ~ model + hard3 + lift_type, cluster SE by task_id
%%   all6: hard3 OR=0.149, p=1.073e-07; Composite OR=0.655, p=0.4858
%%   strong3: hard3 OR=0.092, p=1.945e-07; Composite OR=0.654, p=0.5535
%%   five-no-GLM: hard3 OR=0.151, p=1.456e-06; Composite OR=0.731, p=0.6236
%%   Safe claim: hard3 is the robust difficulty axis; lift type is descriptive only.
%%
%% PASS-CONDITIONED COMPACTNESS (UNPAIRED; descriptive survivor sets only):
%%   Pro   n=115, median RRES=0.993, median copy=0.957
%%   Flash n=108, median RRES=0.998, median copy=0.968
%%   Luna  n=102, median RRES=0.815, median copy=0.164
%%   GLM   n= 68, median RRES=1.013, median copy=0.947
%%   Qwen  n= 63, median RRES=0.975, median copy=0.757
%%   OSS   n= 36, median RRES=0.983, median copy=0.180
%% PAIRED COMPACTNESS (same-task common passes):
%%   Pro vs Luna: n=97; median copy 0.966 vs 0.191; Pro>Luna 76,
%%                Luna>Pro 18, ties 3; Wilcoxon p=1.18e-13;
%%                matched-pairs rank-biserial r=0.881;
%%                median RRES 0.993 vs 0.697.
%%   Flash vs Luna: n=92; median copy 0.970 vs 0.204; Flash>Luna 75,
%%                  Luna>Flash 14, ties 3; Wilcoxon p=4.24e-13;
%%                  matched-pairs rank-biserial r=0.885;
%%                  median RRES 0.999 vs 0.713.
%%   Safe claim: comparable functional success can conceal systematically different
%%   extraction strategies. DO NOT claim lower copying makes Luna globally "better".
%%
%% SEMANTIC L1 (AI-assisted assistant-first-pass; NOT human gold):
%%   Valid Pro+Flash artifact-failure rows n=63 (Pro 28, Flash 35), after task-defect
%%   rows are excluded from this semantic denominator.
%%   behavior_drift             54 = Pro 26 + Flash 28
%%   contract_api_completion     7 = Pro  2 + Flash  5
%%   packaging_modularization    2 = Pro  0 + Flash  2
%%   closure-class total        61/63 (drift + API completion)
%%   localization label = 0 in this L1 slice; DO NOT generalize to "localization solved".
%%
%% PROCESS / TOKEN SECONDARY DIAGNOSTICS:
%%   Empty submissions: Pro 0, Flash 0, Luna 6, GLM 38, Qwen 25, OSS 2.
%%   Qwen empties are tool-validation/security events; process != semantic failure.
%%   Token accounting differs by provider; DO NOT rank all six by token use.
%%   Comparable incremental accounting:
%%     Pro overall median=88,020; P90=225,314; steps median=41.5; P90=80.1
%%     Flash overall median=123,618; P90=269,526; steps median=63.0; P90=116.2
%%     Pro Core vs hard3 median tokens = 81,509 vs 118,318 (p=9.72e-06)
%%     Flash Core vs hard3 median tokens = 108,258 vs 202,895
%%   Safe secondary claim: hard3 elicits more interaction from Pro/Flash while pass
%%   rates still fall sharply; additional interaction alone does not remove hard-tail gap.
%%
%% CRITICAL NON-CLAIMS:
%%   - DO NOT report Qwen as 63/125 or GLM as 68/112; empty submissions are failures.
%%   - DO NOT claim Pro significantly outperforms Flash (p=0.09229).
%%   - DO NOT claim GLM significantly outperforms Qwen (p=0.5682).
%%   - DO NOT claim Direct>Adapted>Composite as an independent difficulty finding.
%%   - DO NOT infer contract closure from Hidden failure alone.
%%   - DO NOT treat semantic L1 as human-gold annotation.
%%   - DO NOT mix official extra 50-task release split into headline Python-150 claims.
%% =====================================================================

%% Central macros for headline numbers used repeatedly in prose.
\newcommand{\HeadlineTasks}{150\xspace}
\newcommand{\HeadlineRepos}{127\xspace}
\newcommand{\ReleaseTasks}{200\xspace}
\newcommand{\ReleaseRepos}{176\xspace}
\newcommand{\BestPassRate}{76.7\%\xspace}
\newcommand{\WorstPassRate}{24.0\%\xspace}
\newcommand{\AllUnsolved}{28\xspace}
\newcommand{\AllUnsolvedHard}{22\xspace}
\newcommand{\IsolationResidualN}{4\xspace}
\newcommand{\SemanticLNN}{63\xspace}
\newcommand{\ClosureLNN}{61\xspace}
\newcommand{\ProLunaCommonN}{97\xspace}
\newcommand{\ProLunaCopyPro}{0.966\xspace}
\newcommand{\ProLunaCopyLuna}{0.191\xspace}
\newcommand{\ProLunaWilcoxon}{1.18\times10^{-13}}
\newcommand{\ProLunaRBC}{0.881\xspace}

\lstset{
  basicstyle=\ttfamily\small,
  keywordstyle=\bfseries,
  commentstyle=\itshape\color{gray},
  breaklines=true,
  frame=single,
}

%% Conference metadata: keep or replace according to the actual target venue.
\setcopyright{acmlicensed}
\copyrightyear{2026}
\acmYear{2026}
\acmConference[FSE '26]{ACM International Conference on the Foundations of Software Engineering}{July 5--9, 2026}{Montreal, Canada}
%% \acmSubmissionID{...}

\begin{document}

\title{FeatureLiftBench: Benchmarking Behavior-Preserving Feature Lifting in Repository-Level Code Agents}

\author{Anonymous Author(s)}

\begin{abstract}
Extracting an existing capability from a mature software repository is different
from repairing that repository in place or implementing a new feature from a
standalone specification. An agent must identify relevant implementation evidence,
recover the APIs, dependencies, state, resources, and edge semantics required for
behavioral completeness, and reconstruct the capability behind a new package
boundary without retaining a runtime dependency on the source project. We
introduce \flb, a benchmark for repository-level, behavior-preserving feature
lifting. The released Python suite contains \ReleaseTasks tasks from \ReleaseRepos pinned repositories;
our headline study uses a frozen \HeadlineTasks-task split under a Full-Repository / No-Hint
protocol. Agents receive the complete repository and a public behavioral contract,
but no source-location hints, benchmark tests, protected evaluator, or reference
implementation. A submission passes only if it builds, satisfies Public and Hidden
behavior derived from the same contract, and remains functional after access to
the source repository is removed. Across six OpenHands model--harness
configurations, Functional Pass@1 ranges from \WorstPassRate to \BestPassRate, and the strongest
configuration still fails 35 of 150 tasks. Failure analysis shows that independence
is rarely the final binding constraint: only \IsolationResidualN delivered artifacts across the
full matrix pass Build, Public, and Hidden but fail Isolation. On the two strongest
backends, an AI-assisted close-read of artifact failures is dominated by incomplete
recovery of required observable behavior, which we characterize as a
\emph{contract-closure gap}. Functional success also hides sharply different
extraction strategies: on \ProLunaCommonN tasks passed by both DeepSeek V4 Pro and GPT-5.6 Luna,
median copied-code fraction is \ProLunaCopyPro versus \ProLunaCopyLuna (paired Wilcoxon
$p=\ProLunaWilcoxon$). These results show that feature lifting requires evaluating
not only whether an agent produces a correct artifact, but also where behavioral
closure fails and how the successful artifact is constructed.
\end{abstract}

\begin{CCSXML}
<ccs2012>
 <concept>
  <concept_id>10011007.10011074.10011099.10011102</concept_id>
  <concept_desc>Software and its engineering~Software testing and debugging</concept_desc>
  <concept_significance>500</concept_significance>
 </concept>
 <concept>
  <concept_id>10011007.10011074.10011111.10011113</concept_id>
  <concept_desc>Software and its engineering~Software verification and validation</concept_desc>
  <concept_significance>300</concept_significance>
 </concept>
</ccs2012>
\end{CCSXML}

\ccsdesc[500]{Software and its engineering~Software testing and debugging}
\ccsdesc[300]{Software and its engineering~Software verification and validation}

\keywords{code agents, repository-level software engineering, benchmark, feature extraction, program modularization, behavioral validation}

\maketitle

%% =====================================================================
\section{Introduction}
%% =====================================================================

Software maintenance frequently requires engineers to extract an existing
capability from a larger system. A parser may need to become a reusable library,
a configuration resolver may need to move into another service, or a plugin
registry may need to be separated from the framework that originally hosted it.
Such work appears in modernization, dependency reduction, service decomposition,
legacy migration, and software reuse. Although the desired capability already
exists, extraction is not a mechanical copy operation. Its behavior may depend
on helper functions, exception hierarchies, stateful registries, packaged
resources, configuration defaults, and third-party libraries distributed across
the repository. Copying the most obvious implementation region can therefore
produce an artifact that succeeds on a happy path while silently changing edge
semantics or retaining hidden dependencies on the source project.

Repository-level code agents appear well suited to this setting. They can search
large workspaces, inspect call sites and tests, edit multiple files, execute
commands, and package new artifacts. Yet common evaluation paradigms ask
different questions. Issue-resolution benchmarks evaluate whether an agent can
repair a repository until its tests pass~\cite{jimenez2024swebench}; code-generation
benchmarks evaluate synthesis from a specification~\cite{chen2021codex}; and
feature-location research evaluates whether relevant implementation regions can
be found~\cite{dit2013featurelocation}. None of these settings directly asks
whether an agent can recover an existing capability from an entangled repository,
preserve its declared behavior, and make the result execute independently of the
source project.

We call this capability \emph{repository-level behavior-preserving feature
lifting}. Feature lifting combines three obligations:
\begin{enumerate}
  \item \textbf{Locate:} identify repository evidence relevant to the requested
        capability;
  \item \textbf{Close:} recover the APIs, helpers, dependencies, state,
        resources, and edge semantics needed for behavioral completeness; and
  \item \textbf{Isolate:} package the recovered capability so that it executes
        independently of the source repository.
\end{enumerate}
This decomposition is an explanatory model, not a causal factorization. A failed
artifact does not by itself reveal which obligation failed; semantic attribution
requires evaluator, artifact, and when available trajectory evidence.

We introduce \flb, a benchmark for this setting. Each task provides a complete
repository pinned to an immutable source revision together with a public
behavioral contract. Under the Full-Repository / No-Hint protocol, the agent may
inspect the entire repository but receives no source-location hints, benchmark
tests, protected evaluator, or reference implementation. It must produce a
standalone \texttt{featurelifted} package. The evaluator checks installation and
imports, primary Public behavior, deeper Hidden behavior derived from the same
public contract, and independence after the original source repository becomes
unavailable.

A benchmark of this kind is only useful if its evaluation obligations are
traceable to information legitimately available to the agent. We therefore pair
construction with a provenance-aware validation pipeline. For each task, Public
and Hidden behavioral cases are linked to public contract clauses and source
repository evidence. An LLM-assisted \validatoragent scales semantic inspection
by identifying missing provenance, ambiguous requirements, underdetermined
behavior, and evaluator inconsistencies. The agent is not the correctness oracle:
final task retention additionally requires deterministic execution, feasible
reference behavior, source isolation, freeze checks, and targeted human review.

The released Python suite contains \ReleaseTasks tasks from \ReleaseRepos pinned repositories. Our
headline empirical study uses a frozen Python-150 split from \HeadlineRepos repositories;
the additional 50-task release split is reported separately in the appendix.
Tasks span Direct, Adapted, and Composite lifts across parsing, configuration,
validation, resource handling, registries, protocols, algorithms, and workflow
behavior. Functional correctness and extraction compactness are evaluated
separately.
\emph{Functional Pass@1} requires all correctness gates to pass, whereas
\emph{Reference-Relative Extraction Size} (\rres) describes the footprint of
functionally passing artifacts relative to an evaluator-side feasible reference.
This separation prevents broad vendoring and incomplete minimal implementations
from being conflated in a single score.

Our empirical study asks four research questions:
\begin{description}
  \item[RQ1 --- Capability.] How well can current coding agents perform
        repository-level feature lifting under a fixed model--harness protocol?
  \item[RQ2 --- Failure.] At which evaluation gates do agents fail, and which
        constraints remain binding for strong systems?
  \item[RQ3 --- Difficulty.] How is empirical task difficulty distributed, and
        which benchmark construction factors remain associated with success after
        controlling for model backend?
  \item[RQ4 --- Extraction quality.] Among functionally successful artifacts, how
        do extraction footprint and copied-code behavior differ across agents on
        the same tasks?
\end{description}

The results reveal a consistent picture. Current systems span a wide capability
range (24.0\%--76.7\% Functional Pass@1) without saturating the benchmark. Strong
backends almost never fail to emit or build a package; their remaining errors are
concentrated at behavioral evaluation. Across the full six-model matrix, only four
delivered artifacts pass Build, Public, and Hidden but fail Isolation, indicating
that independence is rarely the final binding constraint. A close-read of the
Pro+Flash artifact-failure slice further suggests that the residual strong-model
gap is dominated by incomplete recovery of the required behavioral contract.
Finally, functional success alone does not characterize extraction behavior:
DeepSeek V4 Pro and GPT-5.6 Luna have much closer pass rates than their radically
different copied-code fractions on their common passing tasks.

This paper makes three contributions:
\begin{itemize}
  \item \textbf{Task and benchmark.} We formalize repository-level,
        behavior-preserving feature lifting and construct \flb, a 200-task Python
        benchmark from 176 pinned real-world repositories with standalone
        artifact requirements and deterministic evaluation.
  \item \textbf{Construction and validation methodology.} We develop a
        provenance-aware pipeline that connects source evidence, public contracts,
        protected evaluation cases, feasible references, and isolation checks.
        \validatoragent provides scalable semantic auditing while deterministic
        checks and targeted human review safeguard final task validity.
  \item \textbf{Empirical findings.} We benchmark six model backends under a fixed
        OpenHands execution protocol and characterize capability, binding failure
        constraints, structured task difficulty, contract-closure symptoms, and
        pass-conditioned extraction behavior.
\end{itemize}

%% =====================================================================
\section{Related Work}
\label{sec:related-work}
%% =====================================================================

\subsection{Repository-Level Software-Engineering Benchmarks}

Executable software-engineering benchmarks have established the importance of
real repositories, tool use, sandboxed execution, and test-based verification.
SWE-bench asks models or agents to resolve real GitHub issues in existing
repositories~\cite{jimenez2024swebench}. SWE-agent demonstrates that the agent--computer
interface itself can materially affect repository-level performance~\cite{yang2024sweagent}.
Terminal-oriented benchmarks extend executable evaluation to broader command-line
workflows~\cite{merrill2026terminalbench}. These benchmarks typically retain the
original repository as the execution target: success is defined by modifying or
repairing that workspace. \flb instead treats the source repository as evidence
from which a new artifact must be reconstructed. Passing requires that the lifted
package continue to work after the source project is no longer available.

Recent work also emphasizes that executable-agent scores depend on the combined
model--harness configuration rather than on the base model alone. We adopt this
reporting discipline by fixing the OpenHands runtime, prompt, context policy, and
budget in the official \mainprotocol leaderboard and varying only the model
backend. Runtime comparisons, if reported, remain separate from model scores.
\draftnote{Complete the related-work review with the final set of recent
repository-level agent benchmarks and primary-source citations.}

\subsection{Feature Development and Repository-Level Generation}

Repository-level feature development asks an agent to add requested behavior to
an existing codebase. It therefore shares with feature lifting the need to
understand distributed repository context and to produce multi-file software
changes. The output boundary is nevertheless different. Feature-development
systems typically extend the original project, whereas feature lifting must
recover behavior that already exists and reconstruct it behind a new independent
package boundary. The source repository may guide the implementation, but it is
not part of the final runtime environment. This distinction makes dependency,
resource, and package-boundary closure part of correctness rather than merely
implementation style.
\draftnote{Add direct comparison and citations for the most relevant recent
feature-development / repository-generation benchmarks.}

\subsection{Code Generation, Repository Understanding, and Localization}

Code-generation benchmarks evaluate whether a model can synthesize a function or
program from natural-language or executable specifications~\cite{chen2021codex}.
Feature-location research instead studies how developers and tools identify files,
entities, or regions relevant to a concern; Dit et~al. survey classical feature
location and its evidence sources~\cite{dit2013featurelocation}. Feature lifting
requires both inference and artifact construction but is reducible to neither.
Locating the correct implementation is insufficient when transitive behavior,
resources, or public APIs are omitted. Conversely, a behaviorally equivalent
rewrite may succeed even when it copies little upstream code.

\subsection{Program Slicing, Modularization, and Software Extraction}

Feature lifting is related to program slicing~\cite{weiser1984slicing},
multi-dimensional separation of concerns~\cite{tarr1999degrees}, refactoring,
dependency analysis, and library extraction. Classical slicing methods typically
assume a program representation and an explicit slicing criterion. \flb instead
evaluates an autonomous agent given a natural-language behavioral contract and a
complete repository, and grades the delivered independent artifact by observable
behavior. The benchmark does not claim to compute a minimal semantic slice. Its
frozen references are feasible comparison points, and \rres is a compactness proxy
rather than an optimality certificate.

\begin{table*}[t]
  \caption{Conceptual positioning of FeatureLiftBench relative to common
  repository-level evaluation settings.}
  \label{tab:positioning}
  \centering
  \begin{tabularx}{\textwidth}{@{}p{0.20\textwidth}p{0.25\textwidth}p{0.25\textwidth}X@{}}
    \toprule
    \textbf{Setting} & \textbf{Input} & \textbf{Output} & \textbf{Primary objective} \\
    \midrule
    Issue repair & Repository + issue / failing behavior & Modified source repository & Repair existing project \\
    Feature addition & Repository + new requirement & Modified source repository & Add new behavior \\
    Code generation & Specification & New implementation & Synthesize requested behavior \\
    Localization & Repository + query / concern & Files, symbols, or regions & Find relevant implementation \\
    \textbf{Feature lifting} & \textbf{Repository + public contract} & \textbf{Independent artifact} & \textbf{Preserve and isolate existing capability} \\
    \bottomrule
  \end{tabularx}
\end{table*}

%% =====================================================================
\section{The FeatureLiftBench Benchmark}
\label{sec:benchmark}
%% =====================================================================

\subsection{Task Definition}
\label{sec:task-definition}

A task consists of a pinned source repository $R$, a public behavioral contract
$S$, an initial workspace $W$, and a protected evaluator $J$. Given
$(R,S,W)$, a model backend $M$ embedded in a fixed agent harness $H$ produces a
package $P$:
\begin{equation}
  P = \operatorname{Run}(M,H,R,S,W).
  \label{eq:run}
\end{equation}
The official protocol requires the delivered artifact to live under the
submission package boundary and to execute without importing or otherwise
accessing the source project at runtime.

The evaluator assigns gate-level outcomes:
\begin{equation}
  \operatorname{Functional}(P)
  = B(P) \land P_{\mathrm{pub}}(P)
    \land H_{\mathrm{hid}}(P) \land I(P),
  \label{eq:functional}
\end{equation}
where $B$ checks build and importability, $P_{\mathrm{pub}}$ checks primary
observable behavior, $H_{\mathrm{hid}}$ checks deeper combinations and boundary
behavior from the same published contract, and $I$ checks independence from the
source repository. A valid Hidden check must map to a requirement that is
observable from the public contract; an unmapped private requirement is a task
defect rather than a model failure.

The task deliberately permits multiple implementation strategies. Extraction,
adaptation, and behaviorally equivalent reimplementation are legal. The benchmark
does not reward reproducing a particular source slice. Instead, it evaluates
whether the delivered artifact satisfies the public contract after the source
project is removed.

\begin{figure}[t]
  \centering
  \fbox{%
    \parbox{0.92\linewidth}{\centering
      Pinned repository + public contract\\[2pt]
      $\downarrow$\\[2pt]
      Full-Repository / No-Hint agent execution\\[2pt]
      $\downarrow$\\[2pt]
      Standalone \texttt{featurelifted} package + trajectory\\[2pt]
      $\downarrow$\\[2pt]
      Build $\rightarrow$ Public $\rightarrow$ Hidden $\rightarrow$ Isolation\\[2pt]
      $\downarrow$\\[2pt]
      Functional Pass + pass-conditioned compactness}}
  \caption{FeatureLiftBench evaluation pipeline. Protected tests, evaluator code,
  and evaluator-side references are never exposed during agent execution.}
  \Description{A pipeline from a pinned repository and public contract through
  agent execution and a standalone package to build, public, hidden, and isolation
  evaluation, followed by functional and compactness reporting.}
  \label{fig:evaluation-pipeline}
\end{figure}

\subsection{Benchmark Construction Pipeline}
\label{sec:construction}

The benchmark is constructed as a sequence of explicit transformations and
validation gates rather than by asking an LLM to invent a task end to end.
Programmatic checks establish repository identity and executable behavior; LLMs
assist with semantic candidate discovery, contract drafting, and validation; and
human review is reserved for ambiguous or high-risk cases. Figure~\ref{fig:construction-pipeline}
shows the high-level process.

\begin{figure}[t]
  \centering
  \fbox{%
    \parbox{0.92\linewidth}{\centering
      Real-world repository selection\\
      $\downarrow$\\
      Feature candidate discovery\\
      $\downarrow$\\
      Feature boundary and required-API construction\\
      $\downarrow$\\
      Public behavioral contract\\
      $\downarrow$\\
      Feasible reference + Public/Hidden evaluator\\
      $\downarrow$\\
      \validatoragent provenance and ambiguity audit\\
      $\downarrow$\\
      Deterministic build, behavior, isolation, and freeze validation\\
      $\downarrow$\\
      Final benchmark instance}}
  \caption{High-level benchmark construction and validation pipeline. LLM
  assistance is used for semantic work, but executable checks and provenance
  constraints determine benchmark eligibility.}
  \Description{A benchmark construction pipeline from repositories through feature
  discovery, contracts, references and tests, validator-agent auditing, and final
  deterministic validation.}
  \label{fig:construction-pipeline}
\end{figure}

\paragraph{Repository selection and freezing.}
Candidate repositories are acquired at fixed revisions. Each retained task records
its source identity, revision, and digest in a canonical registry. The paper suite
is evaluated only against frozen materializations identified by the active suite
manifest. Source mismatch is treated as an infrastructure or eligibility defect,
not a model outcome.

\paragraph{Feature candidate discovery.}
Candidate capabilities are selected when they correspond to a plausible reuse or
modularization target and expose behavior that can be evaluated independently.
Candidate discovery can use repository structure, symbols, documentation, tests,
resources, and LLM-assisted semantic inspection. A candidate is rejected when the
target is trivial, inseparable from the complete application boundary, dependent
on unavailable infrastructure, or not deterministically testable offline.

\paragraph{Boundary and contract construction.}
For each retained candidate, construction records the required package surface,
behavioral obligations, allowed dependencies, and any preservation requirements.
The public contract describes observable behavior rather than revealing a source
location or a reference extraction. Contracts are designed to support multiple
behaviorally valid implementations while making evaluator obligations recoverable
from the information given to the agent.

\paragraph{Reference and evaluator construction.}
Each task includes an evaluator-side feasible implementation or extraction used to
establish solvability and to provide a compactness reference. Public cases cover
primary contract behavior; Hidden cases stress deeper combinations, boundary
conditions, exception behavior, state transitions, and preservation requirements
from the same contract. References and evaluators are never exposed during model
execution.

\paragraph{Deterministic validation.}
A candidate enters the suite only if the reference installs and passes all
behavioral and isolation checks in the frozen offline environment. Integrity
checks reject tasks that can receive credit through forbidden source imports,
protected-test access, path leakage, undeclared dependency channels, or evaluator
bypass. \draftnote{Insert the exact automated construction-gate counts from the
frozen benchmark manifest.}

\begin{table*}[t]
  \caption{Roles of automation, LLM assistance, and human review in benchmark
  construction. The exact implementation details should match the released
  construction pipeline.}
  \label{tab:construction-roles}
  \centering
  \begin{tabularx}{\textwidth}{@{}p{0.27\textwidth}p{0.18\textwidth}p{0.18\textwidth}X@{}}
    \toprule
    \textbf{Stage} & \textbf{Programmatic} & \textbf{LLM-assisted} & \textbf{Human role} \\
    \midrule
    Repository freeze & Yes & No & Review acquisition exceptions \\
    Candidate discovery & Partial & Yes & Review uncertain candidates \\
    Contract drafting & Partial & Yes & Resolve ambiguity / scope \\
    Evaluator construction & Yes & Optional & Review high-risk cases \\
    Reference validation & Yes & No & Inspect failures \\
    Contract provenance audit & Yes & Yes (\validatoragent) & Adjudicate flagged cases \\
    Isolation and integrity & Yes & No & Inspect exceptional failures \\
    Final eligibility & Yes & Supporting & Targeted adjudication \\
    \bottomrule
  \end{tabularx}
\end{table*}

\subsection{Contract Provenance and Validator-Agent}
\label{sec:validator}

Construct validity is especially important for Hidden evaluation. If a protected
test requires behavior that is neither stated nor reasonably implied by the public
contract, the benchmark measures access to private information rather than feature
lifting. We therefore represent validation as a provenance relation among protected
evaluation obligations, contract clauses, and source evidence.

For a Hidden behavioral case $h_i$, the validator seeks a mapping
\begin{equation}
  h_i \rightarrow s_j \rightarrow e_k,
  \label{eq:provenance}
\end{equation}
where $s_j$ is the public contract clause that licenses the check and $e_k$ is
repository evidence establishing that the behavior belongs to the source
capability. The source-evidence link does not require the agent to copy the source
implementation; it establishes that the benchmark is preserving an existing
feature rather than inventing a new private requirement.

\validatoragent receives the task's pinned repository, public contract, required
API metadata, and evaluator-side behavioral descriptions available to the
validation pipeline. It produces structured audit records rather than a binary
correctness verdict. The audit covers at least the following dimensions:
\begin{enumerate}
  \item \textbf{Requirement provenance:} whether each protected behavioral
        obligation maps to an explicit or defensible public contract clause;
  \item \textbf{Source provenance:} whether the public obligation corresponds to
        behavior supported by the pinned source repository;
  \item \textbf{Ambiguity:} whether multiple materially different interpretations
        of the public requirement would be reasonable;
  \item \textbf{Underdetermination:} whether repository evidence and the public
        contract are insufficient to determine the protected expectation;
  \item \textbf{Evaluator consistency:} whether Public, Hidden, reference, and
        isolation requirements conflict or imply incompatible behavior; and
  \item \textbf{Boundary integrity:} whether a task can pass through protected
        information access or forbidden source dependencies.
\end{enumerate}

The validator is deliberately not the final benchmark oracle. LLM judgments can
be unstable and can share model-family biases with evaluated systems. Validator
outputs therefore serve as scalable evidence collection and triage. Flagged cases
are resolved by deterministic execution and, when semantic adjudication is still
required, targeted human review. We report validator provenance, unknown or
unresolved rates, and the exact retention policy rather than presenting AI-assisted
labels as independent human gold.

\draftnote{Before submission, replace this paragraph with the exact frozen
validator protocol: model(s), prompt/version hash, number of independent passes,
label schema, automatic rejection rules, human-review rule, and aggregate
Explicit/Recoverable/Ambiguous/Underdetermined counts.}

\subsection{Suite Composition and Taxonomy}
\label{sec:suite}

The released suite, Python-200$'$, contains 200 tasks from 176 pinned
repositories (Table~\ref{tab:suite}); the frozen Python-150 split is the headline
evaluation set. It emphasizes Python libraries, developer
tools, and framework components. One source repository may support multiple
distinct tasks when the requested capabilities and evaluation boundaries are
separable.

\begin{table}[t]
  \caption{Composition of the Python-200$'$ paper suite.}
  \label{tab:suite}
  \centering
  \begin{tabular}{lrrl}
    \toprule
    \textbf{Split} & \textbf{Tasks} & \textbf{Repositories} & \textbf{Role} \\
    \midrule
    Frozen Python-150 & 150 & 127 & Headline evaluation \\
    Additional-50 & 50 & 50 & Appendix/release extension \\
    \midrule
    Python-200$'$ & 200 & 176 & Released suite \\
    \bottomrule
  \end{tabular}
\end{table}

Tasks are organized along three descriptive axes. \emph{Feature family} records
the requested behavior, including parsing, serialization, configuration,
validation, resource handling, registry behavior, caching, protocols, algorithms,
and workflows. \emph{Primary entanglement} records the dominant mechanism that
makes independent extraction difficult: data-model, parser-state, framework,
configuration/environment, resource, legacy, or third-party entanglement.
\emph{Lift type} characterizes the relationship between the requested artifact
and the source implementation:
\begin{itemize}
  \item \textbf{Direct:} the target can largely be transferred while preserving
        its original structure;
  \item \textbf{Adapted:} the target requires boundary adaptation, dependency
        replacement, or interface reshaping; and
  \item \textbf{Composite:} the target behavior must be reconstructed from
        multiple source regions or interacting abstractions.
\end{itemize}
These labels support sampling and analysis; they are not shown to the agent and do
not affect Functional Pass.

The combined suite contains 68 Direct, 100 Adapted, and 32 Composite tasks. On
the headline Python-150 split, lift type is treated as a descriptive structural
taxonomy rather than an independent difficulty axis: Direct and Composite are
strongly confounded with the in-suite construction split, and the apparent raw
Direct--Adapted--Composite gradient disappears after controlling for that split.
The additional 50-task release extension broadens repository and behavior
coverage but is not used as the paper's main difficulty claim. \draftnote{Verify the final taxonomy
counts and repository count against the frozen paper manifest.}

An earlier easy expansion exhibited ceiling behavior and copy-heavy solutions.
We retain that result only as benchmark-development evidence motivating stronger
copy resistance and task diversity; it does not contribute to the main
Python-200$'$ leaderboard. Detailed development history is deferred to the
appendix or artifact documentation.

\subsection{Full-Repository / No-Hint Protocol}
\label{sec:protocol}

The official \mainprotocol protocol controls external task and execution
conditions while varying the model backend. The agent receives the complete
source tree and public behavioral contract. It may inspect repository files, run
shell commands, write code, install only permitted dependencies available in the
frozen environment, and execute legal self-tests. It does not receive source
entrypoint hints, benchmark Public tests, Hidden tests, evaluator source, or the
feasible reference.

\begin{table*}[t]
  \caption{Controlled and varying factors in the official \mainprotocol evaluation.}
  \label{tab:factors}
  \centering
  \begin{tabularx}{\textwidth}{@{}p{0.34\textwidth}X@{}}
    \toprule
    \textbf{Factor} & \textbf{Official Main treatment} \\
    \midrule
    Task contract and repository snapshot & Fixed per task \\
    Initial workspace and source registry & Fixed per task \\
    Agent runtime & Fixed to OpenHands \mainprotocol \\
    Prompt and tool interface & Fixed \\
    Action/step budget and context envelope & Fixed \\
    Evaluator and Docker image & Fixed and digest-attested \\
    Model backend & Varied in the main matrix \\
    Public benchmark tests & Withheld in Main \\
    Source-location hint & Withheld in Main \\
    Evaluator, reference, and Hidden tests & Never exposed \\
    \bottomrule
  \end{tabularx}
\end{table*}

Each task attempt produces four evidence layers: (1) the final artifact, (2) the
agent trajectory and process status, (3) token, step, latency, context, and
resource statistics, and (4) deterministic evaluator outcomes. The evaluator
result is authoritative for Functional Pass. Agent completion status is reported
separately because a run can terminate after leaving a passing artifact or declare
completion while its artifact still fails.

\subsection{Metrics}
\label{sec:metrics}

\paragraph{Functional Pass@1.}
A task receives one functional point only if all gates in
Equation~\ref{eq:functional} pass. For diagnosis, functional failures receive one
exclusive first-failure stage:
\begin{equation}
  \text{missing} \rightarrow \text{build} \rightarrow \text{public}
  \rightarrow \text{hidden} \rightarrow \text{isolation}.
\end{equation}
Infrastructure failures and evidence-availability defects are excluded or reported
separately according to the frozen eligibility policy.

\paragraph{Reference-Relative Extraction Size.}
For a functionally passing submission $P$ with evaluator-side feasible reference
$P_{\mathrm{ref}}$, we report
\begin{equation}
  \operatorname{RRES}(P)
  = \frac{\operatorname{normalized\_size}(P)}
         {\operatorname{normalized\_size}(P_{\mathrm{ref}})}.
  \label{eq:rres}
\end{equation}
\rres is accompanied by copied-code, file-count, and dependency-footprint
statistics where available. It is reported only on functional passes. Cross-model
compactness comparisons additionally report same-task subsets because survivor
sets can differ substantially.

\paragraph{Process diagnostics.}
Tokens, steps, latency, completion status, and empty submissions are reported as
secondary diagnostics rather than correctness metrics. Token accounting differs
across providers, so cross-provider token rankings are avoided; same-accounting
comparisons are reported only where they are meaningful.

%% =====================================================================
\section{Experimental Setup}
\label{sec:experimental-setup}

\subsection{Models and Agent Runtime}

The headline study evaluates six model backends on the same frozen Python-150 task
IDs using OpenHands Official Main. The task contract, repository snapshot, initial
workspace, prompt, tool interface, context envelope, step budget, timeout, agent
image, evaluator image, and scoring rule are fixed; only the model backend varies.
Each task is attempted once. Functional Pass is determined by the final collected
artifact and deterministic evaluator rather than by runner completion status.

\begin{table}[t]
  \caption{Model backends in the frozen Python-150 evaluation.}
  \label{tab:models}
  \centering
  \begin{tabular}{ll}
    \toprule
    \textbf{Backend} & \textbf{Serving / profile} \\
    \midrule
    DeepSeek V4 Pro & API backend \\
    DeepSeek V4 Flash & API backend \\
    GPT-5.6 Luna & OpenLux \\
    GLM-5.3-Flash & API backend \\
    Qwen3.6-35B-A3B-FP8 & self-hosted \\
    GPT-OSS 120B & self-hosted / served backend \\
    \bottomrule
  \end{tabular}
\end{table}

\subsection{Execution and Evaluation Environment}

The active freeze is identified by task-manifest digest
\texttt{6c20ff0307762503a73cbb9ff32e9992c6446e4b17483a68373027be58cbf419}.
Runs use a 120-step envelope, 128k context, Full-Repository / No-Hint prompting,
and a 3600-second task timeout. Evaluation occurs in a Dockerized offline
environment with digest-attested agent and evaluator images. The evaluator installs
the collected package, applies Build, Public, Hidden, and Isolation checks, and
records gate-level outcomes. Empty submissions count as failures. Runner
\texttt{status} is not used as the score.

\subsection{Statistical Protocol}

For Functional Pass@1 we report Wilson 95\% confidence intervals. Pairwise model
comparisons use exact McNemar tests over the same 150 task IDs. The task-difficulty
analysis uses descriptive solve-frequency distributions and a logistic model of
pass/fail with model backend, the in-suite \texttt{hard3} construction split, and
lift type as predictors, with standard errors clustered by task. Pass-conditioned
copy-fraction comparisons are restricted to same-task intersections and use paired
Wilcoxon signed-rank tests with matched-pairs rank-biserial effect sizes. Because
each main cell contains one attempt, the results characterize Pass@1 under this
fixed model--harness configuration and do not estimate run-to-run stochastic
variance.

\section{Benchmark Results}
\label{sec:results}

\subsection{RQ1: Overall Feature-Lifting Capability}
\label{sec:rq1}

Table~\ref{tab:main} summarizes the overall FeatureLiftBench results
along three complementary dimensions: functional correctness,
extraction quality, and interaction efficiency.

Functional Pass@1 varies substantially across the six evaluated
model--harness configurations, from 24.0\% to 76.7\%.
DeepSeek V4 Pro achieves the highest observed pass rate
(115/150, 76.7\%), followed by DeepSeek V4 Flash
(108/150, 72.0\%) and GPT-5.6 Luna (102/150, 68.0\%).
Even the strongest configuration fails 35 of the 150 tasks,
indicating that the benchmark remains far from saturated.

The ranking by functional correctness does not translate directly
into a ranking by extraction characteristics.
Among functionally successful artifacts, Pro and Flash have median
copy fractions of 0.957 and 0.968, respectively, whereas Luna has a
median copy fraction of 0.164.
The corresponding mean copy fractions are 0.777, 0.810, and 0.415.
Similarly, the mean RRES values can be substantially larger than their
medians---for example, Pro has mean RRES 1.781 but median RRES 0.993---
revealing a right-tailed distribution in which a minority of successful
artifacts are substantially larger than the feasible reference.

These pass-conditioned aggregate statistics should not be interpreted
as direct paired model comparisons because different models succeed on
different task subsets. We therefore perform same-task paired
compactness analyses in Section~\ref{sec:rq4}.

% ============================================================
% Main FeatureLiftBench leaderboard
% ============================================================
\begin{table*}[t]
  \caption{
    Overall results on the frozen FeatureLiftBench Python-150 suite.
    Functional Pass@1 is computed over all 150 tasks.
    RRES and copy fraction are computed only over functionally passing
    artifacts. Steps and tokens are measured over all 150 runs.
    Lower RRES and copy fraction indicate more compact and less
    source-copy-heavy extractions, respectively, but should not be
    interpreted as semantic superiority by themselves.
  }
  \label{tab:main}
  \centering
  \scriptsize
  \setlength{\tabcolsep}{3.6pt}
  \resizebox{\textwidth}{!}{
  \begin{tabular}{l
                  c
                  ccc
                  ccc
                  cc
                  cc}
    \toprule
    & \multicolumn{1}{c}{\textbf{Correctness}}
    & \multicolumn{3}{c}{\textbf{Extraction Size}}
    & \multicolumn{3}{c}{\textbf{Source Copying}}
    & \multicolumn{2}{c}{\textbf{Steps}}
    & \multicolumn{2}{c}{\textbf{Tokens}} \\
    \cmidrule(lr){2-2}
    \cmidrule(lr){3-5}
    \cmidrule(lr){6-8}
    \cmidrule(lr){9-10}
    \cmidrule(lr){11-12}

    \textbf{Model}
    & \textbf{Pass@1 $\uparrow$}
    & \textbf{Mean RRES}
    & \textbf{Median RRES}
    & \textbf{IQR RRES}
    & \textbf{Mean Copy}
    & \textbf{Median Copy}
    & \textbf{IQR Copy}
    & \textbf{Median}
    & \textbf{P90}
    & \textbf{Median}
    & \textbf{P90} \\
    \midrule

    DeepSeek V4 Pro
    & \textbf{115/150 (76.7\%)}
    & 1.781
    & 0.993
    & [0.769, 1.290]
    & 0.777
    & 0.957
    & [0.716, 0.987]
    & 41.5
    & 80.1
    & 88,020$^{\dagger}$
    & 225,314$^{\dagger}$ \\

    DeepSeek V4 Flash
    & 108/150 (72.0\%)
    & 1.730
    & 0.998
    & [0.893, 1.221]
    & 0.810
    & 0.968
    & [0.839, 0.990]
    & 63.0
    & 116.2
    & 123,618$^{\dagger}$
    & 269,526$^{\dagger}$ \\

    GPT-5.6 Luna
    & 102/150 (68.0\%)
    & 0.953
    & 0.815
    & [0.310, 1.002]
    & 0.415
    & 0.164
    & [0.048, 0.967]
    & 35.5
    & 74.4
    & -- 
    & -- \\

    GLM-5.3-Flash
    & 68/150 (45.3\%)
    & 2.473
    & 1.013
    & [0.917, 1.158]
    & 0.843
    & 0.947
    & [0.811, 0.979]
    & 85.5
    & 127.0
    & --
    & -- \\

    Qwen3.6-35B
    & 63/150 (42.0\%)
    & 1.556
    & 0.975
    & [0.650, 1.228]
    & 0.645
    & 0.757
    & [0.326, 0.972]
    & 41.5
    & 106.6
    & 1,818,620$^{\ddagger}$
    & 5,521,030$^{\ddagger}$ \\

    GPT-OSS 120B
    & 36/150 (24.0\%)
    & 1.176
    & 0.983
    & [0.292, 1.648]
    & 0.305
    & 0.180
    & [0.007, 0.530]
    & 19.0
    & 61.1
    & 525,960$^{\ddagger}$
    & 1,877,320$^{\ddagger}$ \\

    \bottomrule
  \end{tabular}
  }

  \vspace{2pt}
  \begin{minipage}{0.99\textwidth}
    \footnotesize
    \textit{Notes.}
    RRES and copy statistics are conditioned on Functional Pass, with
    $n=115,108,102,68,63,36$ for Pro, Flash, Luna, GLM, Qwen, and OSS,
    respectively.
    $^{\dagger}$ Pro and Flash report incremental tokens
    (uncached prompt tokens + completion tokens).
    $^{\ddagger}$ Qwen and GPT-OSS report provider-reported
    \texttt{total\_tokens}.
    Luna and GLM runs do not contain a reliable token field.
    Token values from different accounting schemes are therefore not
    directly comparable.
  \end{minipage}
\end{table*}

\paragraph{Finding 1.}
Current coding agents exhibit substantial but incomplete feature-lifting capability:
the benchmark separates model--harness configurations over a 52.7-point range and
is not saturated by the strongest evaluated backend.

\subsection{RQ2: Where Do Feature-Lifting Agents Fail?}
\label{sec:rq2}

Table~\ref{tab:independent-gates} decomposes functional failures across
the four evaluator gates. Unlike a first-failure analysis, these gate
outcomes are non-mutually-exclusive: an artifact may simultaneously
violate multiple behavioral or isolation requirements.

Failures concentrate primarily in the behavioral gates.
Across the 829 delivered artifacts, 244 (29.4\%) fail at least one
Public check and 296 (35.7\%) fail at least one Hidden check.
By comparison, 39 artifacts (4.7\%) fail the raw Isolation check.

More importantly, most Isolation failures are not themselves the
binding reason for functional failure. Only four artifacts across all
six backends pass Build, Public, and Hidden while failing Isolation.
Thus, although source independence is an explicit requirement of
FeatureLiftBench, it is rarely the final constraint preventing a
delivered artifact from passing.

\textbf{Finding 2.}
\emph{Feature-lifting failures are dominated by behavioral correctness
rather than package isolation. Across 829 delivered artifacts, only
4 (0.5\%) satisfy all behavioral gates but fail solely because of the
Isolation requirement.}

\begin{table}[t]
  \caption{
    Empirical task-difficulty spectrum on frozen FeatureLiftBench
    Python-150. Difficulty is characterized by the number of the six
    evaluated backends that solve each task, rather than by a
    predefined task label.
  }
  \label{tab:difficulty-spectrum}
  \centering
  \small
  \setlength{\tabcolsep}{7pt}

  \begin{tabular}{crrr}
    \toprule
    \textbf{Solved by}
    & \textbf{Tasks}
    & \textbf{Share}
    & \textbf{Cumulative $\leq k$} \\
    \midrule

    0/6 & 28 & 18.7\% & 28  (18.7\%) \\
    1/6 &  7 &  4.7\% & 35  (23.3\%) \\
    2/6 &  9 &  6.0\% & 44  (29.3\%) \\
    3/6 & 22 & 14.7\% & 66  (44.0\%) \\
    4/6 & 36 & 24.0\% & 102 (68.0\%) \\
    5/6 & 31 & 20.7\% & 133 (88.7\%) \\
    6/6 & 17 & 11.3\% & 150 (100.0\%) \\

    \bottomrule
  \end{tabular}
\end{table}

\paragraph{Finding 2.}
Functional failures concentrate at behavioral gates. Independence is rarely the
binding constraint after Build, Public, and Hidden behavior are satisfied.

\subsection{RQ3: Empirical Task Difficulty}

Task-level solve frequency reveals a broad empirical difficulty spectrum rather
than a uniformly easy or uniformly impossible suite. Twenty-eight tasks are solved
by none of the six backends, whereas 17 are solved by all six. The hard tail is
strongly concentrated in the in-suite \texttt{hard3} split: 22 of the 28 all-fail
tasks are \texttt{hard3}.

\begin{table}[t]
  \caption{Empirical solve-frequency spectrum on Python-150.}
  \label{tab:difficulty-spectrum}
  \centering
  \begin{tabular}{rrrr}
    \toprule
    \textbf{Solved by} & \textbf{Core-100} & \textbf{hard3-50} & \textbf{Total} \\
    \midrule
    0/6 & 6 & 22 & 28 \\
    1/6 & 0 & 7 & 7 \\
    2/6 & 4 & 5 & 9 \\
    3/6 & 17 & 5 & 22 \\
    4/6 & 31 & 5 & 36 \\
    5/6 & 27 & 4 & 31 \\
    6/6 & 15 & 2 & 17 \\
    \bottomrule
  \end{tabular}
\end{table}

The Core--hard3 contrast is large for every backend except GPT-OSS: Pro drops from
94\% to 42\%, Flash from 90\% to 36\%, Luna from 82\% to 40\%, GLM from 62\%
to 12\%, and Qwen from 55\% to 16\%. A raw analysis also appears to show a
Direct--Adapted--Composite gradient. However, lift type is strongly confounded with
construction split: 53/56 Direct tasks are Core, while 15/18 Composite tasks are
\texttt{hard3}. In a logistic model controlling for backend and \texttt{hard3},
Composite is not significantly different from Direct (OR$=0.65$, $p=0.554$),
whereas \texttt{hard3} remains strongly associated with lower success. We therefore
treat lift type as descriptive taxonomy rather than an independent difficulty
finding.

\paragraph{Finding 3.}
The robust difficulty axis in Python-150 is the in-suite \texttt{hard3}
construction split; the apparent lift-type gradient is confounded with split
composition.

\subsection{RQ4: Extraction Quality on Common Successes}

Correctness and compactness are evaluated separately because each model passes a
different subset of tasks. Table~\ref{tab:compactness} first reports descriptive
pass-conditioned summaries. The key comparison is paired: among 97 tasks passed by
both Pro and Luna, median copied-code fraction is 0.966 for Pro and 0.191 for Luna;
Pro copies more on 76 tasks, Luna on 18, with three ties. The paired Wilcoxon test is
highly significant ($p=1.18\times10^{-13}$) with a large matched-pairs
rank-biserial effect ($r=0.881$). Flash versus Luna shows the same pattern on 92
common passes ($p=4.24\times10^{-13}$, $r=0.885$).

\begin{table*}[t]
  \caption{Pass-conditioned extraction statistics on Python-150. Unpaired medians
  are descriptive because survivor sets differ; statistical comparison uses
  same-task intersections.}
  \label{tab:compactness}
  \centering
  \begin{tabular}{lrrr}
    \toprule
    \textbf{Backend} & \textbf{$n$ pass} & \textbf{Median RRES} & \textbf{Median copy fraction} \\
    \midrule
    DeepSeek V4 Pro & 115 & 0.993 & 0.957 \\
    DeepSeek V4 Flash & 108 & 0.998 & 0.968 \\
    GPT-5.6 Luna & 102 & 0.815 & 0.164 \\
    GLM-5.3-Flash & 68 & 1.013 & 0.947 \\
    Qwen3.6-35B & 63 & 0.975 & 0.757 \\
    GPT-OSS 120B & 36 & 0.983 & 0.180 \\
    \bottomrule
  \end{tabular}
\end{table*}

\paragraph{Finding 4.}
Comparable functional success can conceal systematically different extraction
strategies. Functional Pass alone is therefore insufficient to characterize the
artifacts produced by feature-lifting agents.

\section{Diagnostic Analysis}
\label{sec:analysis}

\subsection{From Behavioral Failure to Contract Closure}

Gate outcomes identify where an artifact fails, not why. We therefore perform an
AI-assisted L1 close-read of all valid artifact-level failures from the two
strongest backends (Pro and Flash; $n=63$ after excluding task-defect rows). The
analysis uses the public contract, first-failure evaluator evidence, submitted
package, and trajectory evidence when needed. It is an assistant first pass rather
than independent human gold and is interpreted as diagnostic evidence.

\begin{table}[t]
  \caption{Primary L1 symptom on Pro+Flash artifact-level failures ($n=63$).}
  \label{tab:semantic}
  \centering
  \begin{tabular}{lrrr}
    \toprule
    \textbf{Primary symptom} & \textbf{Total} & \textbf{Pro} & \textbf{Flash} \\
    \midrule
    Behavioral drift & 54 & 26 & 28 \\
    Contract/API completion & 7 & 2 & 5 \\
    Packaging/modularization & 2 & 0 & 2 \\
    \bottomrule
  \end{tabular}
\end{table}

Sixty-one of 63 coded rows fall into behavioral drift or contract/API completion.
In this slice, the typical failure is therefore not the absence of a package but a
package whose observable behavior remains incomplete. We call this recurring
residue the \emph{contract-closure gap}. The term is deliberately narrower than
``Hidden failure'': a closure label requires evidence that the violated obligation
is licensed by the public contract and that the delivered artifact omits or changes
it.

\subsection{Representative Cases}

\paragraph{Alembic revision maps.}
Both Pro and Flash implement \texttt{get\_revision}, but treat the input
\texttt{base} as the symbolic base of the map, shadowing a revision whose literal
identifier is also \texttt{base}. The relevant API is present; the boundary
semantics are wrong.

\paragraph{aiohttp URL parameters.}
The submitted artifact contains exception and token checks, yet a required
invalid-name path still fails to raise. Repository inspection and broad
implementation coverage therefore coexist with an open contract branch.

\paragraph{Copy-heavy successful extraction.}
On \texttt{blinker\_\_signal\_registry\_core\_\_001}, multiple backends pass
all functional gates with substantially different reference-relative footprints.
The case illustrates why a functional pass is an eligibility condition for, rather
than a substitute for, extraction-quality analysis.

\subsection{Secondary Process and Cost Diagnostics}

Process failures are reported separately from artifact-level failures. Qwen's 25
empty submissions are tool-validation/security events; GLM has 38 empty
submissions; Pro and Flash have none. These outcomes are retained as failures in the
fixed model--harness score but are not reinterpreted as semantic contract failures.
Token accounting is not comparable across all providers. Within Pro and Flash,
where the same incremental accounting is available, \texttt{hard3} tasks consume
more tokens than Core while achieving much lower pass rates (Pro median 81.5k vs
118.3k; Flash 108.3k vs 202.9k). This secondary evidence indicates that additional
interaction alone does not eliminate the hard-tail gap.

\section{Discussion}
\label{sec:discussion}

\subsection{Feature Lifting Is a Distinct Agent Capability}

Feature lifting differs from issue repair because the source repository is evidence
rather than the final execution target. An agent must determine which parts of the
source environment are semantically necessary, reconstruct them behind a new
package boundary, and satisfy an output contract after the original repository is
removed. It also differs from localization: finding or copying relevant code can be
necessary while remaining insufficient for complete observable behavior.

\subsection{The Residual Strong-Model Gap Is Behavioral}

The gate analysis and L1 close-read point in the same direction. Strong backends
rarely fail to emit or build a package, and Isolation is seldom the final binding
constraint once behavioral gates pass. Their remaining artifact failures are
instead dominated by missing API obligations or behavioral drift. This motivates
agent designs that track contract obligations and their supporting repository
evidence explicitly, rather than treating additional search, reflection, or test
iterations as sufficient by themselves.

\subsection{Correctness and Extraction Quality Are Different Dimensions}

A broad vendoring solution can preserve behavior while offering little compact
modularization; a compact package can be incomplete. The strong paired difference
between Pro and Luna shows that even similar functional capability can arise from
systematically different extraction strategies. We therefore keep Functional Pass
and RRES/copy statistics separate rather than combining them into a single scalar
leaderboard.

\subsection{What the Difficulty Analysis Does and Does Not Show}

The in-suite \texttt{hard3} split defines a robust empirical hard tail across
medium and strong backends. By contrast, the raw Direct--Adapted--Composite pattern
is confounded with construction split and does not survive the controlled logistic
analysis. We therefore avoid interpreting lift type as a causal or independent
difficulty factor. More generally, task-level taxonomies describe this curated
suite; they are not population estimates over software repositories.

\subsection{Scope of the Claims}

All main scores characterize model backends embedded in the fixed OpenHands
Official Main configuration. They are not measurements of base models independent
of runtime, tool interface, prompt, context policy, provider behavior, or recovery
logic. Process events such as tool validation can therefore affect the official
model--harness score and are reported separately from semantic artifact analysis.

\section{Threats to Validity and Responsible Release}
\label{sec:threats}
%% =====================================================================

\paragraph{Construct validity.}
Functional Pass is a deterministic but finite operationalization of behavioral
completeness. Public and Hidden cases are mapped to one published contract, but
finite tests cannot prove semantic equivalence. A Hidden requirement that is not
licensed by the public contract is a task defect, not an agent failure. We address
this risk through provenance records, \validatoragent auditing, deterministic
reference checks, and targeted human adjudication. We additionally report
sensitivity views that exclude unresolved or conservative audit cases when
appropriate. \rres, copied-code fraction, file count, and dependency footprint are
proxies and can miss semantic copying or penalize legitimate alternative designs.

\paragraph{Internal validity.}
Results depend on source materialization, dependency locks, prompt and context
policy, agent/evaluator images, model endpoints, and runner state. We therefore
separate infrastructure defects from model outcomes and require source,
dependency, context, provenance, and image attestations for the paper bundle.
Development pilots performed under non-identical conditions are not pooled into a
single causal estimate.

\paragraph{External validity.}
Python-200$'$ is a maintainer- and AI-assisted selection of Python libraries,
tooling, and framework components. It is not a random sample of software
repositories and does not establish performance on GUI applications, large product
systems, GPU software, distributed services, or non-Python ecosystems. Popular
upstream projects may also appear in model training data. Repository diversity
does not guarantee balanced behavior or entanglement mechanisms.

\paragraph{Statistical conclusion validity.}
The main leaderboard uses one attempt per task. This design measures Pass@1 under
a fixed harness but does not estimate stochastic run-to-run variance. We report
Wilson confidence intervals and exact paired task outcomes where appropriate, but
these do not replace repeated-run sensitivity. Model APIs and serving stacks can
change despite stable model names; exact run dates and profiles are recorded.

\paragraph{Annotation and validator validity.}
Task taxonomy, provenance records, and semantic failure labels include AI-assisted
curation. LLM-assisted validation can be inconsistent, can share biases with the
models being evaluated, and cannot substitute for an executable oracle. We
therefore disclose validator model and prompt provenance, retain unknown or
unresolved outcomes, and report human inter-rater agreement only when raters are
genuinely independent humans.

\paragraph{Runtime dependence.}
Official Main fixes OpenHands. Scores may change under another agent runtime, tool
interface, context policy, or recovery strategy. Runtime comparisons are reported
separately unless evaluated under a controlled paired protocol.

\paragraph{Licensing and release.}
Upstream repositories retain heterogeneous licenses. Release packaging, source
acquisition, evaluator-side references, and generated submissions must preserve
applicable license obligations. Large source and result archives may be distributed
through checksummed acquisition instructions rather than committed directly.
Protected evaluation assets should be released in a form that preserves benchmark
integrity while supporting independent reproduction of aggregate results.

%% =====================================================================
\section{Conclusion}
\label{sec:conclusion}

We introduced \flb, a benchmark for repository-level, behavior-preserving feature
lifting. The released Python suite contains 200 tasks from 176 pinned repositories,
and our headline frozen Python-150 evaluation compares six model backends under a
fixed Full-Repository / No-Hint OpenHands protocol. Functional Pass separates
Build, Public, Hidden, and Isolation correctness from pass-conditioned extraction
footprint.

Current systems span 24.0\%--76.7\% Functional Pass@1 and do not saturate the
benchmark. Strong backends predominantly fail at behavioral evaluation rather than
package construction, while only four delivered artifacts in the full six-model
matrix pass Build, Public, and Hidden but fail Isolation. An AI-assisted close-read
of Pro+Flash artifact failures further points to a contract-closure gap dominated
by missing API obligations and behavioral drift. Finally, common-success analysis
shows that similar functional capability can conceal radically different extraction
strategies. Together, these results motivate code agents that reason explicitly
about behavioral obligations and support evaluation that distinguishes correctness
from how the reusable artifact is constructed.


%% =====================================================================
%% Optional appendix. Keep only material that is fully frozen and auditable.
%% =====================================================================
\appendix

\section{Artifact and Reproducibility Checklist}
\label{app:reproducibility}

The released paper bundle should include, at minimum, the following immutable
metadata for each official run:
\begin{itemize}
  \item task identifier, source repository revision, and source digest;
  \item public contract version and task manifest digest;
  \item evaluator and container image digests;
  \item model/provider profile, run date, runtime revision, and prompt hash;
  \item context envelope, action/step budget, timeout, and dependency policy;
  \item final package tree, deterministic evaluator output, and trajectory;
  \item token, step, latency, completion, and infrastructure status; and
  \item validator provenance and any human-adjudication record associated with
        task eligibility.
\end{itemize}


\section{Frozen Main Quantitative Record}
\label{app:frozen-record}

For auditability, the source file also contains a non-rendered frozen experimental
record with the exact denominators, paired-test counts, gate counts, difficulty
spectrum, compactness statistics, semantic-L1 denominator, and non-claims used by
the manuscript. The headline freeze is Python-150 under OpenHands Official Main
(one attempt per task; 120 steps; 128k context; Full-Repository / No-Hint). Empty
submissions are functional failures, and runner \texttt{status} is not the score.
The active task-manifest digest is
\texttt{6c20ff0307762503a73cbb9ff32e9992c6446e4b17483a68373027be58cbf419}.

\begin{table*}[t]
  \caption{Non-exclusive gate failures on delivered Python-150 packages. A row can
  fail multiple gates. Isolation residual means Build, Public, and Hidden pass but
  Isolation fails.}
  \label{tab:all-gates}
  \centering
  \begin{tabular}{lrrrrrrr}
    \toprule
    \textbf{Backend} & \textbf{Has pkg.} & \textbf{Build fail} & \textbf{Public fail}
      & \textbf{Hidden fail} & \textbf{Isolation fail} & \textbf{Iso. given Build}
      & \textbf{Iso. residual} \\
    \midrule
    DeepSeek V4 Pro   & 150 & 0  & 25 & 29  & 0  & 0 & 0 \\
    DeepSeek V4 Flash & 150 & 0  & 26 & 36  & 1  & 1 & 1 \\
    GPT-5.6 Luna      & 144 & 3  & 30 & 35  & 3  & 0 & 0 \\
    GLM-5.3-Flash     & 112 & 3  & 34 & 38  & 5  & 2 & 2 \\
    Qwen3.6-35B       & 125 & 6  & 41 & 58  & 8  & 2 & 0 \\
    GPT-OSS 120B      & 148 & 18 & 88 & 100 & 22 & 4 & 1 \\
    \midrule
    Total & 829 & 30 & 244 & 296 & 39 & 9 & 4 \\
    \bottomrule
  \end{tabular}
\end{table*}

\begin{table}[t]
  \caption{Selected exact paired statistics used in the main text.}
  \label{tab:frozen-paired}
  \centering
  \begin{tabular}{lrrr}
    \toprule
    \textbf{Comparison} & \textbf{A-only} & \textbf{B-only} & \textbf{$p$} \\
    \midrule
    Pro vs Flash & 10 & 3 & 0.09229 \\
    Pro vs Luna  & 18 & 5 & 0.01062 \\
    GLM vs Qwen  & 27 & 22 & 0.5682 \\
    \bottomrule
  \end{tabular}
\end{table}

\section{Secondary and Development-Time Diagnostics}
\label{app:secondary}

Public-feedback, earliest-sufficiency, context-compression, repair-loop, and other
development-time experiments are not part of the frozen Python-150 main comparison.
They should be reported only with their original task subsets, evaluator versions,
model backends, and information interfaces. We do not pool them into the main
leaderboard or present them as methods that improve FeatureLiftBench.

Token and step summaries are likewise secondary because provider accounting is not
uniform. Same-accounting comparisons for Pro and Flash may be reported to document
that the \texttt{hard3} slice consumes more interaction while remaining markedly
harder.

\section{Validator-Agent Audit Schema}
\label{app:validator-schema}

For reproducibility, the final artifact should document the structured output
schema used by \validatoragent. A recommended task-level record contains:
\begin{itemize}
  \item the protected behavioral case or assertion identifier;
  \item mapped public contract clause(s);
  \item mapped source evidence and file/symbol provenance;
  \item validation class (e.g., explicit, recoverable, ambiguous,
        underdetermined, or unresolved);
  \item evaluator-consistency and boundary-integrity flags;
  \item validator rationale and confidence metadata;
  \item deterministic validation outcomes; and
  \item final human/adjudication disposition when required.
\end{itemize}
The submitted paper should report only classes and rules actually implemented in
the frozen benchmark pipeline.

\bibliographystyle{ACM-Reference-Format}
\bibliography{references}

\end{document}
